% Generated draft: requires mathematical and source review.
\documentclass[11pt]{article}
\usepackage[margin=25mm]{geometry}
\usepackage{fontspec}
\setmainfont{Latin Modern Roman}
\usepackage{amsmath,amssymb,mathtools,unicode-math}
\setmathfont{Latin Modern Math}
\usepackage{enumitem,needspace,xurl,hyperref}
\setlength{\parindent}{0pt}
\setlength{\parskip}{5pt}
\setlength{\emergencystretch}{3em}
\newcommand{\N}{\mathbb{N}}
\newcommand{\Z}{\mathbb{Z}}
\newcommand{\Q}{\mathbb{Q}}
\newcommand{\R}{\mathbb{R}}
\newcommand{\C}{\mathbb{C}}
\newcommand{\F}{\mathbb{F}}
\newcommand{\A}{\mathbb{A}}
\newcommand{\PP}{\mathbb{P}}
\newcommand{\E}{\mathbb{E}}
\newcommand{\eps}{\varepsilon}
\newcommand{\dd}{\,\mathrm{d}}
\begin{document}
\section*{UnsolvedMath ID: 30006314}
\textbf{Source problem code:} catalogue entry
\subsection{Definitions and mathematical statement}
Let \(F\) be a global field, with adele ring \(\mathbb A_F\). The supplied record concerns a quadratic space \(V\) and an orthogonal extension \(W=V\oplus F\) of dimension one greater. The quadratic forms, including the form on the added line, are not specified. Write
\[
G=\operatorname{SO}(V)\times\operatorname{SO}(W),\qquad H=\operatorname{SO}(V).
\]
The intended period uses the diagonal embedding \(h\mapsto(h,h\oplus 1)\) of \(H\) into \(G\). This interpretation requires compatible nondegenerate quadratic forms; the record does not specify characteristic restrictions or low-dimensional exclusions.

Let
\[
\pi=\pi_V\boxtimes\pi_W\subset\mathcal A_{\mathrm{cusp}}(G)
\]
be a cuspidal automorphic representation that occurs with multiplicity one in the cuspidal spectrum and is tempered at every place. For a decomposable automorphic vector \(\phi=\bigotimes_v\phi_v\), the intended global period is
\[
\mathcal P_H(\phi)=\int_{H(F)\backslash H(\mathbb A_F)}\phi(h,h\oplus1)\,\mathrm dh,
\]
with Tamagawa measure. The intended local period pairing is a matrix-coefficient integral
\[
J_v(\phi_v,\phi_v)=\int_{H(F_v)}
\left\langle\pi_v(h,h\oplus1)\phi_v,\phi_v\right\rangle_v\,\mathrm dh_v,
\]
whenever this integral converges. The source requires the product of the local Haar measures to give the global Tamagawa measure and the product of the local invariant inner products to give the global Petersson inner product. It writes that global inner product as an integral over \(G(F)\backslash G(\mathbb A_F)^1\), with the usual complex conjugation in the second factor. The superscript \(1\) and its compatibility with the period domain require clarification in the exceptional cases where they matter.

\textbf{Intended conjecture, retaining the source's qualification.} The report labels its Conjecture 8 an ``imprecise form'' of the Ichino--Ikeda conjecture. For every such \(\pi\) and every decomposable \(\phi\), it asks for an identity of the form
\[
\left|\mathcal P_H(\phi)\right|^2
=\frac{1}{|S_\pi|}\prod_v^{\prime}J_v(\phi_v,\phi_v).
\]
Here the prime denotes a product normalized using the unramified local period formulas. The supplied record does not define this product completely. It describes \(S_\pi\) as the centralizer of the conjectural global Arthur parameter and also states that Arthur's work gives an unconditional description. Its shorter statement instead refers to a component group. A precise finite group, including any quotient by a center, must be specified before \(|S_\pi|\) has an unambiguous meaning.

The intended expansion of the normalized product has the schematic form
\[
\left|\mathcal P_H(\phi)\right|^2
=\frac{\Delta_G}{|S_\pi|}
\frac{L(\tfrac12,\pi_V\times\pi_W)}{L(1,\pi,\operatorname{Ad})}
\prod_v J_v^{\natural}(\phi_v,\phi_v).
\]
Here \(J_v^{\natural}\) denotes the local pairing after removing its prescribed unramified scalar, and \(\Delta_G\) denotes the group-dependent normalization factor. The source describes that factor as essentially a product of zeta values at positive even integers. Its Proposition 7 attributes to Ichino--Ikeda the unramified local formula
\[
J_v(e_v,e_v)=\Delta_v
\frac{L_v(\tfrac12,\pi_{V,v}\times\pi_{W,v})}{L_v(1,\pi_v,\operatorname{Ad})},
\]
for a spherical vector of norm one and the stated compact-volume normalization. The extraction mixes local and global notation in this passage, so the displayed local subscripts express its intended local reading, which still requires checking. The precise definitions of \(\Delta_v\), \(\Delta_G\), completed versus incomplete \(L\)-functions, and \(J_v^{\natural}\) are absent. Thus the expanded display records the intended identity, rather than a fully specified numerical conjecture.

The source says convergence of the local matrix-coefficient integral needs an argument and mentions strong temperedness of \(H\backslash G\) as a sufficient condition. It also explains that temperedness is a simplifying hypothesis here: the original Ichino--Ikeda formulation is broader, and its genericity assumptions are said to imply expected temperedness through the generalized Ramanujan conjecture. Neither that expectation nor an extension beyond the stated tempered setting is asserted as a theorem in this draft.
\subsection{Short English statement}
For a multiplicity-one tempered cuspidal representation of two adjacent special orthogonal groups, the intended Ichino--Ikeda identity expresses the squared diagonal period through a central Rankin--Selberg value, an adjoint value, normalized local periods, and an Arthur-parameter factor. The supplied report explicitly calls its formulation imprecise; missing quadratic-space hypotheses, normalization conventions, and the exact finite group prevent a fully specified identity.
\subsection{Sources}
\begin{itemize}
\item Local UnsolvedMath registry: permanent ID \(30006314\), source code \texttt{catalogue entry}, title ``The Ichino--Ikeda Period Conjecture'', Oberwolfach Reports Open Problems collection. The supplied record is the evidence used for this draft; repository registry locations are \texttt{data/unsolvedmath/problems.json} and \texttt{catalog/problems.json}.
\item Supplied primary citation: \emph{Arbeitsgemeinschaft: Relative Langlands Duality} (2025), Oberwolfach Reports, pp.\ 811--882, DOI 10.4171/owr/2025/18. The extraction identifies Section 3.4, Proposition 7, and Conjecture 8, attributed to Ichino--Ikeda. \url{https://doi.org/10.4171/owr/2025/18}.
\item Supplied literature source: Beuzart-Plessis, Chaudouard, and Zydor, \emph{Global Gan-Gross-Prasad for unitary groups}. \url{https://pmihes.centre-mersenne.org/articles/10.1007/s10240-021-00129-1/}.
\item Supplied literature source: Boisseau, \emph{Whittaker--Shintani functions and unramified Ichino--Ikeda}. \url{https://doi.org/10.1093/imrn/rnaf271}.
\item Literature assessment dated exactly \texttt{2026-08-22}: the supplied curation labels the problem partially solved and reports proofs for important unitary-group cases and many special cases, including stable and endoscopic cases, while describing the general orthogonal and broader formulation as open. This is supplied dated curation, not independently verified here. The record's \texttt{corrected\_verified} label and its claim that the official report was checked are likewise supplied provenance, not verification performed for this draft.
\end{itemize}
\end{document}
